\documentclass[12pt]{article}

\usepackage{sbc}
\usepackage[utf8]{inputenc}
\usepackage[brazil,english]{babel}
\usepackage{amsmath,amssymb}
\usepackage{booktabs}
\usepackage{hyperref}
\usepackage{microtype}

\title{Título do Artigo: Subtítulo Explicativo}

\author{Primeiro Autor\inst{1}, Segundo Autor\inst{1,2}, Terceiro Autor\inst{2}}

\address{Departamento de Ciência da Computação -- Instituto Federal\\
  \email{\{autor1, autor2\}@instituto.edu.br}
\nextinstitute
  Laboratório de Inteligência Artificial e Pesquisa Aplicada\\
  \email{autor3@lab.org.br}
}

\begin{document}

\maketitle

\begin{abstract}
This meta-paper describes the style to be used in articles for SBC conferences.
It serves as an example for formatting manuscripts with high academic rigor.
\end{abstract}

\begin{resumo}
Este documento apresenta o formato padrão para artigos submetidos aos simpósios e congressos da Sociedade Brasileira de Computação (SBC). O texto aborda a estrutura modular, citações com referências automáticas e formatação de seções.
\end{resumo}

\section{Introdução}
A área de Ciência da Computação e Tecnologia da Informação exige rigor metodológico na condução de experimentos e na redação de artigos científicos.

Trabalhos recentes na área de processamento de documentos e modelos neurais demonstram a relevância da extração estruturada de informação \cite{Chai2024_X7HSG9C2}.

\section{Trabalhos Relacionados}
A literatura acadêmica sobre modelos biomodais e métricas universais de avaliação como ANLS* estabelece parâmetros quantitativos essenciais \cite{Peer2024_W5J79WT6}.

\section{Metodologia}
A metodologia proposta organiza o pipeline em três etapas sequenciais:
\begin{enumerate}
    \item Coleta e higienização dos dados experimentais;
    \item Treinamento e ajuste fino de hiperparâmetros;
    \item Avaliação com validação cruzada k-fold e testes de significância estatística.
\end{enumerate}

\section{Resultados e Discussão}
Os resultados empíricos demonstraram ganhos consistentes de precisão e recall.

\begin{table}[ht]
\centering
\caption{Comparação de Desempenho dos Modelos Avaliados}
\label{tab:resultados}
\begin{tabular}{lcccc}
\toprule
\textbf{Modelo} & \textbf{Precision} & \textbf{Recall} & \textbf{F1-Score} & \textbf{Tempo (ms)} \\
\midrule
Baseline (Regras) & 0.742 & 0.685 & 0.712 & \textbf{12} \\
BERTimbau Base    & 0.885 & 0.862 & 0.873 & 85 \\
DocFusion (0.28B) & \textbf{0.924} & \textbf{0.910} & \textbf{0.917} & 45 \\
\bottomrule
\end{tabular}
\end{table}

\section{Conclusões}
Este trabalho apresentou uma abordagem estruturada para a pesquisa científica em computação, integrando automação de referências e compilação de alta fidelidade.

\bibliographystyle{apalike}
\bibliography{references}

\end{document}
